% BEGIN VOL07-EXPANSION VII17
\section{Supplementary worked examples}
\label{sec:vii17-pedagogy}

\begin{example}[Sphere curvature]\label{ex:vii17-ped-01}
For a sphere of radius \(R\), both principal curvatures have magnitude \(1/R\).  Consequently
\[
K=\frac1{R^2},
\]
while the sign of \(H\) depends on the chosen normal and sign convention for the shape operator.  Gaussian curvature is orientation-independent.
\end{example}

\begin{example}[Cylinder curvature]\label{ex:vii17-ped-02}
A cylinder of radius \(R\) has principal curvatures of magnitudes \(1/R\) and \(0\).  Hence \(K=0\) although the surface is visibly bent in space.  Zero Gaussian curvature is therefore weaker than being an affine plane.
\end{example}

\begin{example}[Saddle point]\label{ex:vii17-ped-03}
At the origin of \(z=x^2-y^2\), the principal curvatures have opposite signs (up to a common orientation sign).  Their product is negative, so the origin is hyperbolic.  The surface bends upward in one principal direction and downward in the other.
\end{example}

\section{Graded supplementary exercises}
The following exercises add a deliberate progression from concrete calculation to structural proof, hypothesis testing, application, and synthesis.

\subsection*{Standard computations and constructions}

\begin{exercise}[Curvature from eigenvalues]\label{exr:vii17-ped-01}
If the principal curvatures are \(2\) and \(3\), compute \(K\) and the arithmetic mean curvature \(H=(k_1+k_2)/2\).
\end{exercise}
\begin{hint}
Use product and average.
\end{hint}
\begin{solution}
\(K=6\) and \(H=5/2\).
\end{solution}

\begin{exercise}[Cylinder invariants]\label{exr:vii17-ped-02}
For principal curvatures \(1/R\) and \(0\), compute \(K\) and \(H\).
\end{exercise}
\begin{hint}
Use the definitions.
\end{hint}
\begin{solution}
\(K=0\) and \(H=1/(2R)\), up to the common sign convention for \(H\).
\end{solution}

\begin{exercise}[Sphere invariants]\label{exr:vii17-ped-03}
For principal curvatures \(1/R,1/R\), compute \(K\) and \(H\).
\end{exercise}
\begin{hint}
Multiply and average.
\end{hint}
\begin{solution}
\(K=1/R^2\) and \(H=1/R\), with \(H\)'s sign depending on orientation convention.
\end{solution}

\begin{exercise}[Hyperbolic point]\label{exr:vii17-ped-04}
Classify a point with \(K<0\).
\end{exercise}
\begin{hint}
Look at the signs of principal curvatures.
\end{hint}
\begin{solution}
The principal curvatures have opposite signs, so the point is hyperbolic.
\end{solution}

\begin{exercise}[Umbilic test]\label{exr:vii17-ped-05}
Classify a point where \(k_1=k_2\).
\end{exercise}
\begin{hint}
Use the definition.
\end{hint}
\begin{solution}
It is an umbilic point.
\end{solution}

\subsection*{Proofs}

\begin{exercise}[Gaussian curvature is orientation-independent]\label{exr:vii17-ped-06}
Prove \(K\) does not change when the unit normal is reversed.
\end{exercise}
\begin{hint}
Both principal curvatures change sign.
\end{hint}
\begin{solution}
The product \((-k_1)(-k_2)=k_1k_2\) is unchanged.
\end{solution}

\begin{exercise}[Mean curvature changes sign]\label{exr:vii17-ped-07}
Show \(H\) changes sign under normal reversal.
\end{exercise}
\begin{hint}
Both principal curvatures change sign.
\end{hint}
\begin{solution}
\((-k_1-k_2)/2=-H\).
\end{solution}

\begin{exercise}[Extrema of normal curvature]\label{exr:vii17-ped-08}
Using Euler's formula, show principal curvatures are the extrema of normal curvature.
\end{exercise}
\begin{hint}
Write \(k_n(\theta)=k_1\cos^2\theta+k_2\sin^2\theta\).
\end{hint}
\begin{solution}
The expression is a convex combination of \(k_1,k_2\), attaining its extreme values in the corresponding principal directions.
\end{solution}

\begin{exercise}[\(K=\det S\)]\label{exr:vii17-ped-09}
Prove Gaussian curvature equals the determinant of the shape operator.
\end{exercise}
\begin{hint}
Diagonalize the self-adjoint operator.
\end{hint}
\begin{solution}
In an orthonormal principal basis, \(S=\mathrm{diag}(k_1,k_2)\), so \(\det S=k_1k_2=K\).
\end{solution}

\subsection*{Counterexamples and hypothesis tests}

\begin{exercise}[\(K=0\) means plane]\label{exr:vii17-ped-10}
Test on a cylinder.
\end{exercise}
\begin{hint}
The cylinder has one zero and one nonzero principal curvature.
\end{hint}
\begin{solution}
False.  A cylinder has \(K=0\) but is not planar.
\end{solution}

\begin{exercise}[\(H=0\) means \(K=0\)]\label{exr:vii17-ped-11}
Test when principal curvatures are \(a,-a\).
\end{exercise}
\begin{hint}
Compute the average and product.
\end{hint}
\begin{solution}
False.  Then \(H=0\) but \(K=-a^2<0\) unless \(a=0\).
\end{solution}

\begin{exercise}[Positive \(K\) forces both principal curvatures positive]\label{exr:vii17-ped-12}
Remember the orientation can reverse both signs.
\end{exercise}
\begin{hint}
Positive product allows both positive or both negative.
\end{hint}
\begin{solution}
False as a signed statement; \(K>0\) means the curvatures have the same sign, while their common sign depends on the chosen normal.
\end{solution}

\subsection*{Applications and investigations}

\begin{exercise}[Curvature classification]\label{exr:vii17-ped-13}
Explain the geometric distinction among elliptic, hyperbolic, and parabolic points.
\end{exercise}
\begin{hint}
Use the sign of \(K\).
\end{hint}
\begin{solution}
At elliptic points \(K>0\) bending has the same sign in principal directions; at hyperbolic points \(K<0\) signs oppose; at parabolic points \(K=0\) at least one principal curvature vanishes.
\end{solution}

\begin{exercise}[Curvature-aware meshing]\label{exr:vii17-ped-14}
Why might a mesh need smaller elements where \(|k_1|\) or \(|k_2|\) is large?
\end{exercise}
\begin{hint}
Large curvature means normals and tangent planes vary rapidly.
\end{hint}
\begin{solution}
Finer elements reduce geometric approximation error in regions of strong bending.
\end{solution}

\subsection*{Challenge problems}

\begin{exercise}[Recover principal curvatures from \(H,K\)]\label{exr:vii17-ped-15}
Show \(k_1,k_2\) are roots of \(\lambda^2-2H\lambda+K=0\).
\end{exercise}
\begin{hint}
Use their sum and product.
\end{hint}
\begin{solution}
Since \(k_1+k_2=2H\) and \(k_1k_2=K\), the monic polynomial with those roots is \(\lambda^2-2H\lambda+K\).
\end{solution}

\begin{exercise}[Umbilic from discriminant]\label{exr:vii17-ped-16}
Show a point is umbilic exactly when \(H^2-K=0\).
\end{exercise}
\begin{hint}
Compute the discriminant of the principal-curvature polynomial.
\end{hint}
\begin{solution}
The discriminant is \(4(H^2-K)=(k_1-k_2)^2\), so it vanishes exactly when \(k_1=k_2\).
\end{solution}

% END VOL07-EXPANSION VII17
